\href{three-generations-healthcare-it_AIIM.pdf}{PDF}{\textbar{}}{\textbar{}}\href{figures.html}{Figures}

\textsuperscript{1}School of Computer Science, Faculty of Sciences, Holon Institute of Technology (HIT), 52 Golomb St., Holon 5810201, Israel

\textsuperscript{2}Intelligent Systems, Afeka Academic College of Engineering, 218 Bnei Efraim St., Tel-Aviv 6910717, Israel

Correspondence: Yehudit Aperstein, apersteiny@afeka.ac.il. Co-author: A. Apartsin, alexanderap@hit.ac.il. ORCID: Apartsin 0009-0000-7007-3529, Aperstein 0000-0001-6390-9463. Submitted to \emph{Artificial Intelligence in Medicine} as a Position Paper.

{Abstract}

Much of what should happen next is decided in conversation and prose, not structured data. When a clinician writes ``repeat the complete blood count in two weeks,'' the intended action and its timing remain in narrative text, and the accountable owner goes unstated. Leaving intent in free text is implicated in failed test-result follow-up, unclosed referral loops, and avoidable diagnostic harm. We take the position that recovering patient-specific intended actions from unstructured clinical communication, in a form a system can schedule, monitor, and close, is a distinct and tractable problem for artificial intelligence in medicine. It goes beyond clinical language processing, which scores fields one at a time, and complements interoperability standards, which represent intent once structured. We organize health information technology into three layers by the unit each makes computable: the record, the clinical fact, and the intended action. We define the Actionable Clinical Record (ACR): a source-grounded representation of one intended action whose target state is executable readiness, upstream of FHIR workflow resources, not a new interchange format. Recovery has two sides: extracting each action, and reasoning over the resulting set through scheduling, contradiction detection, and consistency checking. Because temporal semantics must be computed and traceability kept source-grounded, not generated, we argue for a hybrid neuro-symbolic architecture with output invariants, and evaluation on executable correctness rather than text overlap. A worked example and two separately reported companion studies support record-level tractability. The ACR and its evaluation framework are offered for the community to build on, evaluate, or refute.

\textbf{Keywords:} clinical intent; actionable clinical record; clinical natural language processing; neuro-symbolic reasoning; temporal reasoning; FHIR; workflow automation

\section{Introduction}\label{introduction}

A clinician writes ``repeat the complete blood count in two weeks,'' or tells a patient to return if symptoms recur. Everyone can read the instruction, but its action, timing, any gating condition, and the responsible owner exist only as narrative or go unstated, outside any structured data a system could act on. Recovering such patient-specific intended actions from unstructured clinical communication, in a form a system can schedule, monitor, and close, is an open problem for artificial intelligence in medicine, and the subject of this position paper.

The cost of leaving intent in free text is measurable. Test-result follow-up often fails \cite{ref1}, in one large health system only about one-third of specialist referrals reached a completed visit \cite{ref2}, and such loop-closure failures, rooted in ambiguous responsibility and care-coordination breakdowns \cite{ref3}, contribute to diagnostic error and avoidable harm {[}4-7{]}. These failures are multi-factorial, involving scheduling, access, and care coordination as well as communication; the factor an information system can address is that the intended action often never becomes machine-actionable.

We take a position: computable clinical intent is a distinct computational target. It goes beyond the per-item extraction formulations of clinical natural language processing, which score targets independently, and it is complementary to the interoperability standards that represent intent only once it has been structured. We organize health information technology not by the technologies it adopts but by the unit of information it makes computable. In our reading of the field's documented history \cite{ref8, ref9}, that unit has been the record and then the clinical fact; we argue the next is patient-specific intent. Recovery has two sides: extracting each intended action, and reasoning over the set of actions an encounter produces.

This position paper analyzes the scientific challenges of that layer. We contribute the layer model and its markers (Section 2); the Actionable Clinical Record (ACR) as the layer's atomic object, made concrete with a worked example (Sections 3--5); an analysis of the open inference sub-problems and the reliability contract they impose (Sections 5.3--6); an executable-correctness evaluation framework, with feasibility evidence from two separately reported companion studies (Section 7); and a research agenda with implications and falsifiable predictions (Section 8).

\section{A Model of Computational Layers}\label{a-model-of-computational-layers}

We distinguish a computational layer from an incremental trend by five markers: a new atomic \emph{unit of computability}; an \emph{external driver}; an \emph{enabling technology}; a new \emph{class of computation} the unit permits; and a \emph{residual limitation} that motivates the next layer's object. The layers are \emph{cumulative}, each building on those beneath it (Figure 1).

The layers differ in standing: the first two are retrospective, drawn from a well-documented history \cite{ref8, ref9}, while the third is a \emph{prospective hypothesis}. The layers organize computational objects rather than assert a strict succession: order entry and care plans already make selected patient-specific intent computable, and the frontier is the large residual that remains in communication. The model is complementary to the data--information--knowledge hierarchy \cite{ref10} and digital-maturity models such as EMRAM \cite{ref11, ref12}, which measure adopted capability; our concern is \emph{which} object becomes computable (Table 1).

\begin{figure}[tbp]
\centering
\includegraphics[width=\linewidth]{figures/Figure_1.pdf}
\caption{\textbf{Figure 1.} The three computational layers of health information technology, organized by the unit of information each makes computable. Layers accumulate rather than replace; layer 3 is proposed. Each layer's label names the sense in which its unit is computable: a record is \emph{available} to computers, a fact \emph{interpretable}, and an intended action \emph{operational}. \emph{Alt text:} three stacked boxes, record at the base, clinical state above it, and clinical intent at the top, with an upward arrow indicating that each layer builds on the one below.}
\end{figure}

\begin{table*}[t]
\centering
\caption{\textbf{Table 1.} The three computational layers against these markers. Layers 1 and 2 are retrospective; layer 3 is proposed. For layers 1 and 2 the residual limitation motivates the unit of the next; for the proposed layer 3 it names the open frontier, safe and reliable execution.}
\small\setlength{\tabcolsep}{4pt}%
\setlength{\tabcolsep}{4pt}
\begin{tabularx}{\textwidth}{>{\hsize=0.653\hsize\raggedright\arraybackslash}X >{\hsize=0.881\hsize\raggedright\arraybackslash}X >{\hsize=1.279\hsize\raggedright\arraybackslash}X >{\hsize=1.187\hsize\raggedright\arraybackslash}X}
\toprule
Axis & Layer 1: Record & Layer 2: Clinical state & Layer 3: Clinical intent (proposed) \\
\midrule
Unit of computability & Clinical record & Discrete clinical fact & Intended clinical action (ACR) \\
Question answered & What was documented? & What is known about the patient? & What is supposed to happen? \\
External driver & Adoption policy and reimbursement & Interoperability and quality measurement & Care continuity and patient safety \\
Enabling technology & Transactional databases; messaging & Terminologies; structured data; interoperability; clinical NLP & Language models; symbolic reasoning; workflow semantics \\
Class of computation & Store, retrieve, exchange & Query, reason, predict & Recover, reconcile, schedule, close \\
Canonical output & Document & Coded fact & Actionable Clinical Record \\
Residual limitation & Meaning remains in narrative & Intent remains in communication & Safe, reliable execution (open frontier) \\
\bottomrule
\end{tabularx}
\end{table*}

\section{The First Two Layers: Record and Clinical State}\label{the-first-two-layers-record-and-clinical-state}

The first layer digitized the clinical record and its transactions, replacing the paper chart with electronic capture, storage, retrieval, transmission, and operations such as order entry and results routing. In the United States, HITECH and Meaningful Use drove rapid nationwide EHR adoption \cite{ref13, ref14}, extending the older problem-oriented-record ideal \cite{ref15}. Its enabling technology, the transactional database, stored and moved information without making its clinical meaning computable across systems.

The second layer turned the record into computable clinical \emph{state}, whose atomic unit of computability is the discrete clinical fact. Its enabling stack was broad, not a single technique: structured data entry, controlled terminologies \cite{ref16}, data warehouses, and interoperability standards \cite{ref17, ref18} mattered as much as language processing (MedLEE, cTAKES) \cite{ref19, ref20}. These enabled the first broad wave of computation \emph{on} clinical data: decision support, quality measurement, registries, and research \cite{ref21}. A tension remained: encoding meaning into structured fields captures facts, while much clinical reasoning and many intended actions remain expressed only in narrative \cite{ref22}.

Neither layer is uniformly realized in routine care. Data quality, semantic normalization, and cross-system interoperability remain incomplete, so even layer-2 computability is partial in practice rather than settled. This matters for the argument that follows in two ways: the intent layer does not presuppose that the layers beneath it are finished, and its prospective status is a difference of kind, not merely of maturity. Intent is not under-coded state awaiting better terminologies; it is a different object, produced in communication and lost when only its concepts are indexed.

\section{The Remaining Computational Gap}\label{the-remaining-computational-gap}

The first two layers made clinical \emph{state} computable, but many patient-specific \emph{intended actions} remain non-computable while they are expressed only in unstructured communication. Computerized order entry is cross-layer infrastructure: its transactional storage belongs to the record layer, while the structured request it captures already makes selected patient-specific intent computable. The residual this layer targets is intent that is communicated but never ordered, the discharge recommendations, conditional advice, patient instructions, and handoffs that drive the loop-closure failures of Section 1 \cite{ref1, ref6}.

Three paradigms already represent process, each once someone has structured it. FHIR represents patient-specific intent through \emph{CarePlan}, \emph{ServiceRequest}, and \emph{Task}, and reusable protocols through \emph{PlanDefinition} \cite{ref23}. Computer-interpretable guidelines encode executable protocols {[}24-26{]}. Process mining reconstructs pathways from event logs \cite{ref27, ref28}. All three are complementary to the layer proposed here and often supply the infrastructure that runs its outputs.

The gap is isolated by distinguishing three kinds of process information (Table 2): \emph{normative} (what should generally happen), the domain of guidelines and decision support; \emph{observed} (what actually happened), the domain of records and process mining; and \emph{intended} (what is supposed to happen for a specific patient), often expressed in communication rather than structured directly.

\begin{table*}[t]
\centering
\caption{\textbf{Table 2.} Three kinds of process information. The categories describe the semantic status of a process statement, not exclusive source types. The proposed third layer targets recovery of patient-specific intended process that remains expressed in clinical communication.}
\small\setlength{\tabcolsep}{4pt}%
\setlength{\tabcolsep}{4pt}
\begin{tabularx}{\textwidth}{>{\hsize=0.543\hsize\raggedright\arraybackslash}X >{\hsize=1.179\hsize\raggedright\arraybackslash}X >{\hsize=0.718\hsize\raggedright\arraybackslash}X >{\hsize=1.560\hsize\raggedright\arraybackslash}X}
\toprule
Process type & Meaning & Typical source & Existing paradigm \\
\midrule
Normative & What should generally happen & guideline / protocol & computer-interpretable guidelines; decision support \\
Observed & What actually happened & event log / EHR & process mining \\
Intended & What is supposed to happen for a specific patient & clinical communication & represented once structured (FHIR); recovery from communication proposed here \\
\bottomrule
\end{tabularx}
\end{table*}

The current trajectory of interoperability narrows but does not close this gap. Standardized read APIs are now widely deployed: US Core profiles over the USCDI data classes are required of certified EHRs under United States interoperability and information-blocking rules \cite{ref29}, and patient- and population-level FHIR query is increasingly deployed. Adoption is uneven across the resource space, however. In real-world FHIR applications, use concentrates in read-oriented USCDI resources such as \emph{Patient}, \emph{Observation}, \emph{Condition}, and \emph{Medication}, while the workflow resources that would carry an intended action, \emph{Task}, \emph{ServiceRequest}, and \emph{CarePlan}, are far less commonly reported \cite{ref30}, and by design hold intent only once it has been structured into them \cite{ref23}. The intent itself still originates in unstructured communication, which remains an important carrier of what a clinician means to happen next \cite{ref22}. Recovering that intent, rather than assuming it arrives pre-structured, is the computational problem this layer names.

\section{Computable Clinical Intent}\label{computable-clinical-intent}

The third layer makes patient-specific clinical intent computable by recovering it from communication as a structured representation intended for executable use. Five terms recur: clinical \emph{communication} expresses \emph{intent}; an intent decomposes into intended \emph{actions}; each action is represented by an \emph{Actionable Clinical Record}; a set of related ACRs is the intended \emph{workflow}, a plan; and the sequence actually enacted is the \emph{care process}, the observed process of Table 2.

\subsection{The Actionable Clinical Record}\label{the-actionable-clinical-record}

{Definition 1 (Actionable Clinical Record).} An Actionable Clinical Record is the \emph{recovery target} of the third layer: the FHIR-aligned representation of one patient-specific intended action recovered from communication, carrying the attributes needed to interpret, execute, monitor, or audit it. It is the tuple\\
{ACR = ⟨ action, target, request-intent, actor, temporal~constraint, condition, dependency, status, provenance, confidence ⟩}\\
in which \emph{action}, \emph{request-intent}, \emph{status}, and \emph{provenance} are required (\emph{request-intent} is the authority of the request, a subset of FHIR \emph{RequestIntent}: proposal, plan, order, or option, since ``consider,'' ``recommend,'' and ``order'' a test do not share executable semantics); \emph{target}, \emph{actor}, \emph{temporal~constraint}, \emph{condition}, and \emph{dependency} are populated when communicated or inferred; and \emph{confidence} is a per-attribute map of calibrated recovery-system uncertainty, paired with an \emph{origin} flag marking each value as extracted verbatim, computed deterministically, or inferred. The ACR is not a new interchange format but a projection of FHIR workflow resources onto what must be recovered: a sufficiently populated ACR carries most of the semantics needed to construct a draft FHIR request, once patient context and profile-specific mandatory fields are supplied.

Each attribute has a fixed semantics (Figure 2) and is defined as its FHIR element (Table 3). Two of the definitions carry weight for what follows. \emph{actor} is a role-tagged set that separates the requester, the accountable owner for loop closure, and the intended performer, none of which need be the speaker. \emph{temporal constraint} is a normalized time (absolute, relative, event-anchored, or recurring) rather than a date alone; \emph{dependency} is typed, a relation to an anchor that may be another ACR, a clinical event, an encounter, a result, or an unresolved future trigger, carrying an offset and a resolution state; and \emph{provenance} identifies the source communication, the encounter, the exact span, and, when available, the speaker or author.

\begin{figure}[tbp]
\centering
\includegraphics[width=\linewidth]{figures/Figure_2.pdf}
\caption{\textbf{Figure 2.} The ACR schema. \emph{Alt text:} ten labeled boxes forming the ACR tuple; action, request-intent, status, and provenance are shaded to mark them as required; target, actor, temporal constraint, condition, and dependency are outlined as optional content, and confidence, with its paired origin flag, is per-attribute metadata the system always carries.}
\end{figure}

\begin{table*}[t]
\centering
\caption{\textbf{Table 3.} Each ACR attribute mapped to its FHIR R5 element. The representation is almost entirely available in FHIR, including conditional and dependency structure through \emph{RequestOrchestration} (named \emph{RequestGroup} in FHIR R4, which the companion benchmark uses); a sufficiently populated ACR therefore maps onto a draft FHIR request once patient context and profile-specific mandatory fields are supplied.}
\small\setlength{\tabcolsep}{4pt}%
\setlength{\tabcolsep}{4pt}
\begin{tabularx}{\textwidth}{>{\hsize=0.568\hsize\raggedright\arraybackslash}X >{\hsize=1.531\hsize\raggedright\arraybackslash}X >{\hsize=0.901\hsize\raggedright\arraybackslash}X}
\toprule
ACR attribute & FHIR R5 element & Coverage \\
\midrule
action & \emph{ServiceRequest.code} / \emph{Task.code} & represented \\
target & \emph{ServiceRequest.code} or \emph{bodySite} & represented \\
request-intent & \emph{ServiceRequest.intent} (proposal, plan, order, option: a subset) & represented \\
actor (roles) & \emph{ServiceRequest.requester}, \emph{performer}; \emph{Task.owner} for accountability & represented; \emph{ServiceRequest} has no owner \\
temporal constraint & \emph{occurrence{[}x{]}} (\emph{Timing}/\emph{Period}); \emph{Task.restriction.period} & represented when resolvable \\
condition & \emph{RequestOrchestration.action.condition}; \emph{asNeeded{[}x{]}} for PRN execution & represented \\
dependency & \emph{RequestOrchestration.action.relatedAction}; \emph{basedOn}; \emph{partOf}; \emph{replaces} & represented \\
status (pair) & workflow-status $\rightarrow$ \emph{.status} (draft, active, completed, revoked); recovery-state is ACR-internal & workflow-status represented \\
provenance & \emph{Provenance} resource & resource level; span needs extension \\
confidence & no native element & recovery extension \\
\bottomrule
\end{tabularx}
\end{table*}

The ACR carries four things this infrastructure does not, and they are annotations of the recovery process rather than a rival data model: source-grounding at the level of the exact text \emph{span} rather than the whole resource; per-attribute \emph{calibrated confidence}; an origin flag separating \emph{extracted}, deterministically \emph{computed}, and \emph{inferred} values; and the ability to hold an \emph{unresolved} constraint, a ``before the next visit'' with no scheduled visit, which does not map cleanly to a resolved FHIR occurrence time without a profile extension. These are what make an ACR a scorable recovery target rather than only a message.

Framing the ACR as a target, not a model, is what makes the third layer a well-posed problem: recovery accuracy is undefined against FHIR in general, but it is exactly defined against this schema, and the schema is the unit over which the set-level reasoning of Section 5.3 and the executable-correctness evaluation of Section 7 are computed. The conceptual parallel is executable semantic parsing, as in text-to-SQL, where a fixed target makes recovery measurable by what the structure implies when run rather than by surface overlap. Running an ACR set means resolving the plan it implies against a patient timeline, its enabled and disabled actions, due intervals, owners, ordering, and any detected inconsistency, which the set-level measures of Section 7 score. The contribution here is the recovery target and the executable-correctness evaluation framework, not a new interchange format or a parser.

The third layer spans two distinct axes that together prevent over- and under-claiming. A representational \emph{readiness} ladder runs \emph{mentioned $\rightarrow$ interpreted $\rightarrow$ actionable $\rightarrow$ executable}: recognizing ``complete blood count'' in a note is layer 2; extracting ``repeat complete blood count'' reaches \emph{interpreted} (a typed action and target); adding a normalized temporal constraint or gating condition makes it \emph{actionable}; and resolving an accountable owner and a due date bound to a workflow object makes it \emph{executable}. A partially populated ACR sits below the executable rung until its owner and timing are resolved. Separately, the ACR's \emph{status} is a pair, ⟨~recovery-state, workflow-status~⟩: the recovery-state runs \emph{recovered} then \emph{confirmed} on review, and the workflow-status is the FHIR request lifecycle (\emph{draft}, \emph{active}, \emph{completed}, \emph{revoked}), the two coexisting rather than replacing one another, without changing its representational readiness; this lifecycle maps onto the FHIR request status (recovered and confirmed as \emph{draft}, then \emph{active}, and \emph{completed} or \emph{revoked}), and detecting completion or escalation closes the loop.

\subsubsection{Worked example: from a discharge instruction to closed-loop follow-up}\label{worked-example-from-a-discharge-instruction-to-closed-loop-follow-up}

\textbf{Communication.} Discharge summary, written by the discharging inpatient team and addressed to the primary-care physician (PCP): ``Repeat CBC in two weeks; if palpitations recur, refer to cardiology.''

\textbf{Recovered ACR 1.}\\
{action: repeat-lab} {target: CBC (LOINC 58410-2)} {request-intent: plan} {actor: performer=lab, owner unresolved (candidate PCP)} {temporal: discharge date + 14 d} {condition: none} {status: recovery=recovered, workflow=draft} {provenance: span ``Repeat CBC in two weeks,'' discharge summary} {origin: action/target extracted, temporal computed, owner inferred} {confidence: temporal 0.86, owner 0.55} {(derived) readiness: actionable}

\textbf{Recovered ACR 2.}\\
{action: refer} {target: cardiology} {request-intent: plan} {actor: owner=PCP (inferred)} {condition: palpitations recur (patient-reported)} {temporal: event-anchored (on trigger)} {status: recovery=recovered, workflow=draft} {provenance: span ``if palpitations recur, refer to cardiology''} {origin: action/target/condition extracted, actor inferred} {confidence: owner 0.60}

\textbf{Mapping.} ACR 1 $\rightarrow$ a draft \emph{ServiceRequest} (\emph{intent}~=~plan, \emph{status}~=~draft, \emph{occurrence}~=~discharge~+~14~d, \emph{requester} the discharging team) paired with a follow-up \emph{Task} whose \emph{owner} is left unassigned until the accountable party is confirmed (a \emph{ServiceRequest} carries requester and performer, not an owner); ACR 2 $\rightarrow$ a \emph{RequestOrchestration} action gated on the patient-reported symptom (an applicability condition), a conditional order that comes into existence when the trigger fires, not a PRN flag on an already-active order.

\textbf{Loop closure.} The open \emph{Task} for ACR 1 enters a results queue and is tracked to a resulting CBC; if none is resulted within its +14~d window plus a grace period, the \emph{Task} escalates rather than lapsing silently. ACR 2 stays dormant until its patient-reported trigger fires.

\textbf{Where recovery is hard.} The discharge summary does not say who orders the repeat CBC, so its accountable owner is left unresolved (either the discharging team or the PCP could own it) and the record stays at the \emph{actionable} rung, not \emph{executable}, until an owner is confirmed; the referral's actor is instead inferred from the addressee, at lower confidence. The referral is gated on a patient-observable symptom, so its trigger needs a patient-reported or clinician-monitored channel. A system that silently assigned an owner or a hard date would manufacture certainty the source does not contain.

\textbf{What the layers below cannot supply.} Layer 2 methods can identify the concepts, events, and temporal expressions here; what they do not produce is a single source-grounded operational action that binds them to \emph{repeat}, to a discharge-anchored due window, to a responsible owner, and to the condition that gates the referral. Those bindings, and the flag that the owner is unresolved, are the ACR.

\subsection{From Communication to Workflow}\label{from-communication-to-workflow}

The third layer is a bridge, not a replacement EHR (Figure 3). Its inputs are the range of clinical communication; its output is a set of ACRs that maps onto infrastructure the field already has \cite{ref17, ref23}. The mapping situates the ACR \emph{upstream} of these resources, not beside them; Table 3 gives the element-level correspondence. Dependency and status need profile-specific mapping, because their FHIR relationship and lifecycle semantics differ from the ACR's. Recovery \emph{provenance} can itself be represented with a FHIR \emph{Provenance} resource or profile, while \emph{confidence} generally remains metadata of the recovery system.

FHIR can group related requests (through \emph{RequestOrchestration}), but it holds the result of structuring, not a procedure for recovering it from communication or for checking that a recovered set is mutually consistent, contradiction-free, and feasibly schedulable; the decision-support layers that could check such properties act on structured, already-authored inputs whose workflow semantics are explicit. So the computation this layer names is twofold: recovering each intended action from communication, and reasoning over the resulting set.

\begin{figure}[tbp]
\centering
\includegraphics[width=\linewidth]{figures/Figure_3.pdf}
\caption{\textbf{Figure 3.} ACRs as the bridge between clinical communication and existing workflow infrastructure. \emph{Alt text:} communication sources on the left feed a central clinical-intent layer containing ACR representations, which maps to existing infrastructure such as FHIR resources, schedulers, and referral and results systems on the right.}
\end{figure}

Follow-up is the best-documented member of a broader class the ACR represents: conditional escalation, medication transitions, handoffs, referrals, and pending-result responsibility. Three near-term uses follow. An ambient documentation system can emit ACRs, and hence draft \emph{Task} and \emph{ServiceRequest} entries, rather than prose alone, moving ambient AI from note generation toward order generation. A results-management inbox can carry the recovered follow-up intent alongside each pending result, so that responsibility is explicit rather than assumed. A care-transition handoff can be represented as a set of ACRs whose owners and due times are tracked across the transition rather than restated in free text. Recent shared tasks extract structured orders directly from clinical conversations \cite{ref31}, while ambient documentation shows these conversations are already captured at scale \cite{ref32}.

\subsection{The Scientific Challenges}\label{the-scientific-challenges}

Two mature lines can be mistaken for this layer, and separating them sharpens the open problem. Clinical NLP extracts concepts, events, and temporality {[}19,20,33-38{]}, and, closest to this layer, physician action items and medical decisions from notes \cite{ref39, ref40}. Even order extraction from conversation targets a flat schema scored by text overlap \cite{ref31}; the distinction the ACR draws is operational, requiring that a recovered action be source-grounded, temporally normalized, and carry the actor, condition, and dependency structure needed to \emph{execute} rather than only index it. Computer-interpretable guidelines encode population-level protocols (Section 4); this layer recovers what was instructed for one patient. Table 4 places the ACR against the closest of these extraction schemas.

\begin{table}[tbp]
\centering
\caption{\textbf{Table 4.} The ACR recovery target against the closest clinical-extraction schemas. Each contributes part of the record; none requires a normalized time, an accountable actor, conditional or dependency structure, or set-level coherence, and each is scored by item overlap rather than by executable correctness. CLIP \cite{ref39} tags discharge sentences with action-item types; MedDec \cite{ref40} extracts decision spans; MEDIQA-OE \cite{ref31} extracts conversation orders with utterance-level provenance. A ``no'' marks a field not required by that task or schema, not one the paradigm cannot represent.}
\small\setlength{\tabcolsep}{4pt}%
\small
\setlength{\tabcolsep}{4pt}
\begin{tabularx}{\linewidth}{>{\hsize=1.444\hsize\raggedright\arraybackslash}X >{\hsize=0.833\hsize\raggedright\arraybackslash}X >{\hsize=0.795\hsize\raggedright\arraybackslash}X >{\hsize=0.795\hsize\raggedright\arraybackslash}X >{\hsize=1.133\hsize\raggedright\arraybackslash}X}
\toprule
Capability & CLIP & MedDec & MEDIQA-OE & ACR (target, this work) \\
\midrule
Output unit & sentence label & decision span & order & executable action \\
Typed action & 7 types & categories & order types & yes \\
Source grounding & sentence & span & utterance & span \\
Normalized temporal constraint & no & no & no & yes \\
Accountable actor vs speaker & no & no & no & yes \\
Conditional / dependency structure & no & no & no & yes \\
Set-level coherence & no & no & no & specified (§5.3) \\
Evaluation target & label F1 & span F1 & order overlap & executable correctness \\
\bottomrule
\end{tabularx}
\end{table}

Recovering an executable ACR poses two families of open problems for AI in medicine. The first concerns each intended action in isolation. \emph{Action and argument extraction} must recover not just that a test was mentioned but that it is to be repeated, by whom, and on what object {[}39-41{]}. \emph{Temporal normalization} must turn ``in two weeks'' or ``before the next visit'' into a resolvable constraint. Existing taggers normalize the time expression itself {[}42-45{]}; the harder parts are \emph{anchor selection} (two weeks from admission, discharge, or dictation), which clinical temporal schemes handle with an explicit document-time anchor \cite{ref46}, and constraints that stay unresolved until an event is scheduled, where general-purpose language models remain weak \cite{ref47}. \emph{Actor and responsibility disambiguation} must separate who spoke from who is accountable for acting, which goes beyond speaker-role identification \cite{ref48}: the ordering clinician, the performer, and the narrator can be three different people.

The second family, which distinguishes this layer from both per-item extraction and from workflow representation, concerns the \emph{set} of intents recovered from an encounter. A single communication typically expresses several intents that must be \emph{individuated} rather than returned as a bag of spans. \emph{Conditioning} must attach each conditional action to the proposition that gates it, its observer, and its scope. \emph{Scheduling and ordering} must place actions in time and in dependency order so the result is a feasible plan, not a list of independent timestamps. \emph{Contradiction detection} must flag intents that cannot jointly hold, including a later communication that revises or revokes an earlier one, and \emph{consistency checking} must confirm the set has no orphan dependency or unsatisfiable constraint. \emph{Cross-message reconciliation} is the temporal form of the same requirement, across communications over time, and it builds on coreference and entity linking across records \cite{ref49}. The compared extraction task formulations, which score items independently by text overlap \cite{ref31}, and interoperability standards, which hold structured results rather than recover them, do not provide this reasoning over a recovered set; it is the computation that turns extracted actions into a coherent, executable plan. Across both families, \emph{calibrated abstention} is a cross-cutting requirement: an uncertain recovery is routed to a clinician rather than emitted as an unsafe action. Each requirement has a measure in Table 5, so progress on it is a number.

These set-level relations can be made precise. Two ACRs \emph{contradict} when they act on the same target with mutually exclusive actions, under a stated exclusion relation on action types (hold against continue, cancel against perform), or when they share an action and target with temporal constraints that cannot both hold for a single-occurrence intent. A repeat CBC at one week and again at two weeks are two serial orders unless a later communication supersedes the earlier; that supersession is a \emph{revision}, in which the later ACR replaces the earlier on the same action and target, mapped to \emph{ServiceRequest.replaces} with the earlier status set to \emph{revoked}. Precedence, for resolving which of two contradicting ACRs stands, runs by communication time, then speaker authority, and otherwise the pair is routed to review. A hold with a stated release event is a revision of the held action's temporal constraint; a hold with no release event is a contradiction under the exclusion relation. \emph{Scheduling} treats unconditional event-anchored constraints as a simple temporal problem, whose consistency is decidable in polynomial time \cite{ref50}; conditional actions turn this into a conditional temporal problem, whose consistency is checked per branch of the condition \cite{ref51}. This is the reasoning that clinical guideline and care-plan formalisms already perform over curated inputs, from knowledge-based temporal abstraction and the Asbru time-oriented guideline language \cite{ref52, ref53} to task-network models such as PROforma and GLIF \cite{ref54, ref55}. What is new is performing it over a \emph{recovered}, uncertain, and partially unresolved set, where the reasoner must tolerate abstained attributes rather than assume a complete, consistent plan. The companion evidence of Section 7 grounds the per-record measures; the set-level measures remain to be instantiated on recovered plans, which the evaluation framework specifies and future benchmarks should populate.

\section{Reliability and Architecture}\label{reliability-and-architecture}

A more capable generator does not remove the third layer's requirements: auditability, verifiable temporal correctness, and calibrated deferral are properties of the output contract, not of model capability. Large language models show clinical information-extraction and medical-reasoning capability \cite{ref37, ref56}, with continued gains in recent clinical information-extraction evaluations \cite{ref57}, but can produce unreliable or unsupported outputs \cite{ref58}; this matters most in binding an action to its time, a setting where general-purpose temporal reasoning is known to be weak \cite{ref47}.

These requirements motivate a hybrid design: learned components interpret, explicit components enforce verifiable structure \cite{ref59, ref60}, for example by constraining a generator to a valid schema through grammar-constrained decoding \cite{ref61}. The split follows the semantics: learned extraction proposes candidate actions and arguments \cite{ref41}; deterministic reasoning, over formally specified workflow semantics, resolves times and due dates {[}42-45{]}; and selective prediction routes uncertain cases to a clinician \cite{ref62}, with calibrated confidence so that deferral means something \cite{ref63}. The interface between the components is a typed intermediate representation: the learned component emits a normalized temporal expression together with an explicit anchor label, and the deterministic component computes the due date from it; when the anchor is ambiguous the deterministic step would otherwise inherit a wrong anchor with full confidence, so that case is routed to abstention rather than resolved. Because verbalized and token-level confidence estimates are poorly calibrated for language models, a held-out calibration set supports an abstention rule with finite-sample guarantees on a pre-specified loss under exchangeability: conformal calibration for coverage \cite{ref64}, and conformal risk control for the rate of accepted-but-erroneous records \cite{ref65}. The requirements are invariants: every acted-upon record traces to a source span, speaker, and encounter \cite{ref48}; deterministic semantics are computed, not generated; the system may abstain; higher-consequence actions receive tighter human review \cite{ref66}; an immutable trace links communication to action for audit; and extracted content stays distinct from inferred. Human confirmation is the default until task-specific calibration and prospective safety evidence justify higher automation.

A system that drafts orders from ambient conversation sits on a regulatory and ethical boundary that the output contract is designed to hold. Software a clinician cannot independently review may constitute device software; in particular, time-critical recommendation functions generally do not qualify for the Non-Device CDS exclusion in the United States Clinical Decision Support Software guidance \cite{ref67}; the European Union Artificial Intelligence Act treats AI that is a medical device, or a safety component of one, subject to third-party conformity assessment as high-risk, with obligations for logging, human oversight, and transparency \cite{ref68}. The invariants above, source-linked traceability, computed temporal semantics, abstention, and consequence-scaled review, provide technical mechanisms aligned with those obligations. Two duties are specific to recovery from communication: consent and notice for capturing and processing ambient encounter audio, and a defined failure mode when a recovered conditional action, such as a symptom-gated referral, depends on a patient-observed event that no one is assigned to watch.

\section{Evaluation Framework and Feasibility}\label{evaluation-framework-and-feasibility}

Because third-layer outputs can drive clinical action, they should be evaluated on \emph{executable correctness} rather than text overlap. We specify an evaluation framework for ACR recovery whose measures are reported separately at two levels: at the level of each recovered action, and at the level of the recovered \emph{set}, where a plan's coherence lives (Table 5). The set-level measures are what a per-item extraction score cannot capture, and they are the reason executable correctness is not reducible to token or slot overlap. These assess recovery; loop-closure outcome, a downstream endpoint shaped by workflow rather than recovery correctness, is reported separately.

Every record-level measure presupposes an \emph{alignment} between predicted and reference ACRs, which we compute by optimal bipartite matching under a fixed cost (two points for a matching action, one for a matching coded target, source-span overlap breaking exact ties) that admits a pair only when its score is at least two, below which a prediction is left unmatched. A prediction whose action agrees but whose target differs is matched, with the target scored as an attribute error, so action detection and argument accuracy stay independent. A matched pair is scored on its attributes; an unmatched prediction is a false positive, and an emitted actionable record with no span that expresses it is an unsupported recovery, unsafe only if it crosses the autonomous-execution threshold; an unmatched reference is an omission. Optional attributes are scored where the reference populates them. The scoring penalizes invention, not restraint: an attribute neither grounded in the source nor licensed by a stated inference rule, populated anyway, is an unsupported-attribute error that counts toward an unsafe output; an attribute left unresolved where the reference leaves it unresolved is correct.

\begin{table*}[t]
\centering
\caption{\textbf{Table 5.} The executable-correctness evaluation framework. Per-record measures score each recovered action after alignment; set-level measures score the coherence of the recovered plan. Loop-closure outcome is reported separately as a downstream workflow endpoint.}
\small\setlength{\tabcolsep}{4pt}%
\setlength{\tabcolsep}{4pt}
\begin{tabularx}{\textwidth}{>{\hsize=0.579\hsize\raggedright\arraybackslash}X >{\hsize=1.911\hsize\raggedright\arraybackslash}X >{\hsize=0.510\hsize\raggedright\arraybackslash}X}
\toprule
Measure & Definition (after alignment) & Unit \\
\midrule
\multicolumn{3}{l}{%
Per recovered action} \\
Action detection & precision, recall, and F1 of matched actions & F1 \\
Argument extraction & per-slot F1 for target, actor, temporal, condition, and dependency on matched records & F1 \\
Action--argument linking & accuracy that each argument binds to the correct action & accuracy \\
Temporal-normalization error & difference between predicted and reference resolved time; boundary error for a window; exact match for a symbolic constraint & days / exact \\
Whole-record exact match & fraction of matched records correct on every required and every reference-populated attribute & proportion \\
Unsupported-action rate & fraction of emitted actionable records with no supporting source span & rate (lower) \\
Omitted-action rate & fraction of reference actions with no matched prediction & rate (lower) \\
Calibration / selective risk & record-error risk versus coverage after abstention; expected calibration error of confidence & curve / ECE \\
Provenance accuracy & fraction of records whose provenance span overlaps the reference above a set threshold (for example, token Jaccard 0.5) & rate \\
\multicolumn{3}{l}{%
Per recovered set (plan coherence)} \\
Intent individuation & split and merge rates per reference intent: a reference intent returned as several records, or several returned as one & rate \\
Constraint and dependency recovery & precision, recall, and F1 of the reference's timing, ordering, and dependency constraints, over relation, anchor, offset, and resolution state & P / R / F1 \\
Contradiction and revision detection & F1 of flagged contradictory or revising pairs against the reference & F1 \\
Schedule agreement & share of reference action pairs with a stated or inferred order whose relative order the predicted plan reproduces, evaluated per feasible branch of the conditions, unrecovered pairs counted as disagreement; plan satisfiability reported separately & rate \\
\bottomrule
\end{tabularx}
\end{table*}

An output is \emph{operationally unsafe} when a record crosses the autonomous-execution threshold while a required executable attribute (a resolved time, an accountable owner, or a supporting source span) is missing or wrong, or when the recovered set carries an undetected contradiction. A record deliberately held below that threshold and routed to review is not unsafe for having abstained. Errors are weighted by consequence class, assigned from the action type by a safety taxonomy prespecified for each deployment, so a mis-timed high-consequence order counts for more than a mislabeled routine one, with the weighting reported alongside the unweighted scores. Reference ACRs are annotated with explicit rules for unstated attributes, an actor left unresolved rather than guessed and a temporal anchor chosen by a stated precedence (an explicit anchor, then the discharge or encounter date, then document time), and per-attribute inter-annotator agreement is reported so the attainable ceiling on each measure is visible.

\subsubsection{Worked scoring: a plausible-looking but unsafe recovery}\label{worked-scoring-a-plausible-looking-but-unsafe-recovery}

\textbf{Reference.} The two ACRs of the earlier example: repeat CBC at discharge~+~14~d (owner unresolved); refer to cardiology if palpitations recur.

\textbf{Prediction.} one matched record with wrong bindings, one hallucinated order, one omission:\\
{P1: repeat CBC, discharge + 7 d, actor cardiology} {P2: order chest X-ray (unsupported)} {refer-cardiology: omitted}

\textbf{Alignment.} P1 matches reference ACR 1 (same action and target); P2 is unmatched and unsupported, and, emitted above the autonomous-execution threshold, also operationally unsafe; the referral is an omission.

\textbf{Scores.} action detection precision 1/2, recall 1/2; temporal-normalization error 7~d; actor wrong (cardiology invented for an unstated owner, an unsupported-attribute error); whole-record exact match 0; unsupported-action rate 1/2; omitted-action rate 1/2; the mis-timed, mis-actored order above threshold registers as operationally unsafe.

\textbf{Set-level (a second message).} A later note, ``hold the CBC until after the cardiology visit,'' adds a dependency (CBC after the cardiology visit, an event that is not itself an ACR in the set) and revises ACR 1; the system re-emits P1 unchanged and recovers no dependency. Against the reference: constraint and dependency recovery 0 of 1, the revision is missed (0 of 1), and schedule agreement 0 of 1 (the one reference ordered pair, the visit before the CBC, is not reproduced). The reference plan itself carries a consistency flag, since the CBC is now gated on a cardiology visit that occurs only if palpitations recur, so in the other branch the CBC is never due; a set-level checker routes this to review, while every per-record measure is blind to it.

\textbf{Why it matters.} A text-overlap metric scores P1 highly because it says ``repeat CBC.'' Executable correctness does not: the time is wrong by a week, the accountable owner is invented, an unsupported order is emitted, and a real referral is dropped. The framework surfaces the failures that matter clinically.

Two separately reported companion studies give the framework record-level feasibility evidence, and both isolate the same gap; they are external results \cite{ref69, ref70}, not empirical work of this position paper. A benchmark study realizes the ACR as a FHIR-aligned schema, the Clinical Intent Representation, and harmonizes five existing corpora (CLIP, MedDec, ap\_parsing, PaniniQA, SIMORD/ACI-Bench) into 10,011 intents with a deterministic FHIR R4 mapper \cite{ref69}. Its human-validated core agrees with human decisions on the five critical fields (action type plus the four closed fields) 88.4\% of the time in the audited high-agreement stratum. Evaluated on it at full coverage, five current models given the gold span recover the action type reliably (85 to 91\%) but recover all four closed fields, target, timing, condition, and FHIR request-intent, correctly only 18 to 35\% of the time, and oracle context raises this only to 38 to 47\%. A second study isolates the action--time pairing on a controlled 2,000-note synthetic outpatient corpus: general-purpose models recognize the action almost perfectly (F1 up to 0.99) yet pair it with the correct normalized date only about 0.51 to 0.57 of the time, while a hybrid neural-symbolic pipeline that extracts spans and computes dates deterministically reaches 0.99 Pair F1 with zero-day error, its confidence intervals not overlapping the baselines \cite{ref70}. The two results, from different data and metrics, agree: recognizing an action type from its span was substantially easier than recovering an executable one, and separating learned interpretation from computed structure eliminated date error in that controlled setting once the temporal anchor was fixed, though anchor selection under ambiguity remains open. In these studies, executable correctness, not type accuracy, was the dominant constraint. Neither study reports the set-level measures: the first amounts to whole-record exact match on the closed fields, the second to temporal-normalization error and a per-pair linking F1.

\section{Research Agenda and Implications}\label{research-agenda-and-implications}

The layer model and the ACR define a program in four domains. \textbf{Representation:} a consensus ACR profile aligned to FHIR workflow resources, with attribute semantics, constraint types, span-level provenance, and reconciliation when communications revise or revoke a plan. \textbf{Inference and reliability:} ACR recovery from multilingual and multimodal communication, temporal reasoning, cross-message reconciliation, calibrated abstention, distinguishing patient from clinician intent (intent origin, read from the provenance speaker together with the actor), and separating who spoke from who is responsible for acting. \textbf{Evaluation:} benchmarks on real clinical data, the framework of Section 7 across tasks, and cross-institution generalization. \textbf{Integration and impact:} ACR-to-FHIR write-back into scheduling, results-management, and referral systems \cite{ref23}, review-queue design, regulatory classification, and whether closed-loop operation improves outcomes.

For the AI-in-medicine research community, computable clinical intent supplies a well-posed target and a shared, executable-correctness evaluation framework that separates operationally safe recoveries from merely fluent ones. For standards, the ACR identifies recovery metadata and workflow semantics that could inform implementation guidance or profiles around existing FHIR workflow resources, without introducing a competing interchange format. For practice and policy, closing communication-to-action loops is exactly the kind of measurable safety outcome that quality measurement and interoperability initiatives already reward, and it names where ambient AI can move next, from documenting the encounter to acting on it.

Moving from the synthetic and harmonized corpora of the companion studies to de-identified real-world notes is the central empirical risk. Safe deployment presupposes the invariants of Section 6. The proposal makes testable predictions: evaluation should shift toward executable-correctness measures; reliable systems should expose source-linked, auditable records with selective review; and ambient documentation should extend from notes to tasks and orders \cite{ref32, ref71, ref72}. More decisively, the proposal makes two separately falsifiable claims. First, intent recovery is a distinct task; it is not, if conventional item extraction already captures the provenance, responsibility, temporal, conditional, and uncertainty semantics the ACR requires. Second, set-level reasoning is worthwhile; it is not, if per-item extraction plus direct mapping into existing workflow standards yields plans as coherent, on the set-level measures of Table 5 (individuation, constraint and dependency recovery, schedule agreement), as a method that reasons over the recovered set. Failure of the second claim leaves the recovery task standing; only failure of the first would retire the layer. Separately, the evaluation framework is redundant if standard extraction metrics discriminate operationally unsafe outputs as well as the executable-correctness measures across task classes.

\section{Conclusion}\label{conclusion}

Health IT has made the record and then the clinical fact computable; we take the position that patient-specific clinical intent is the next layer, and that recovering it from communication is a distinct, tractable challenge for artificial intelligence in medicine. What characterizes it is not generating clinical text but recovering, from clinical communication, source-grounded verifiable representations of intended action, formalized as the Actionable Clinical Record upstream of existing workflow infrastructure. We offer the ACR and an executable-correctness evaluation framework, specified sufficiently to instantiate, test, extend, and refute, with task-specific safety thresholds and set-level benchmark construction left to future work.

\section{Data Availability}\label{data-availability}

No new data were generated for this article. The feasibility results in Section 7 come from two companion studies \cite{ref69, ref70}, which report their own data provenance (a synthetic outpatient corpus in one, existing corpora harmonized under their own terms in the other) and availability. The ACR schema (Section 5) and the evaluation framework (Section 7) are specified in the text as an implementable framework; deployment-specific parameters, including execution thresholds, the safety taxonomy, exclusion relations, and benchmark thresholds, are prespecified per task. No custom code was produced.

\section{Author Contributions (CRediT)}\label{author-contributions-credit}

A. Apartsin: Conceptualization, Methodology, Writing -- original draft. Y. Aperstein: Conceptualization, Supervision, Writing -- review and editing.

\section{Funding}\label{funding}

None declared.

\section{Conflicts of Interest}\label{conflicts-of-interest}

None declared.

\section{References}\label{references}
